% !TEX root = ../main.tex

\subsection{Exploratory data analysis}

The training set contains 4{,}958{,}347 rows distributed over 199{,}795 unique search queries; the test set contains 4{,}959{,}183 rows over 199{,}549 searches. Every search returns between 5 and 38 candidate hotels with a sharp mode at 32, which corresponds to the default search-results page size on the Expedia website (see Fig.~\ref{fig:group-sizes}). Both train and test cover the period 2012-11-01 to 2013-06-30, so the partition is a random sample of search queries rather than a temporal split. This simplifies the use of historical priors, because no train search lies temporally before any test search; the priors aggregate the entire train slice without any look-ahead risk.

\begin{figure}[t]
\centering
\includegraphics[width=0.95\linewidth]{eda_group_sizes.png}
\caption{Distribution of candidates per search. Left, histogram with the median at 29 hotels and the 95th percentile at 34. Right, empirical CDF; the distribution is roughly uniform between 5 and 29 candidates and then concentrates sharply around the default page size of 32.}
\label{fig:group-sizes}
\end{figure}

The relevance label is constructed as $\mathrm{rel} = 5 \cdot \mathbf{1}_{\mathrm{book}} + 1 \cdot \mathbf{1}_{\mathrm{click\,only}}$, giving label counts $\{0: 4{,}736{,}468,\ 1: 83{,}489,\ 5: 138{,}390\}$ on the training set. About 95.5\% of rows carry no positive signal at all, which is typical for click-through ranking; a model must learn within-query contrasts rather than pointwise probabilities. The conditional rate $P(\text{book} \mid \text{click}) = 0.62$ shows that the click and book events are tightly coupled in the data and that treating booking and click as two separate sub-targets would be redundant.

Missingness clusters in three groups (Fig.~\ref{fig:missing}). The eight competitor blocks \texttt{comp1\_*} to \texttt{comp8\_*} are 83\% to 98\% missing because they are only filled when a competing online travel agent had a quote available. The two visitor-history columns (\texttt{visitor\_hist\_starrating} and \texttt{visitor\_hist\_adr\_usd}) are 95\% missing, reflecting first-time Expedia visitors. The query-affinity score \texttt{srch\_query\_affinity\_score} and the origin-distance column \texttt{orig\_destination\_distance} are 90\% and 32\% missing respectively. Crucially, LightGBM splits on missingness as a separate value rather than treating NaN as zero, so the missingness pattern itself becomes signal during training; we therefore preserve NaNs rather than impute them.

\begin{figure}[t]
\centering
\includegraphics[width=0.95\linewidth]{eda_missingness.png}
\caption{Per-column missingness rate, sorted within four groups (raw numeric, visitor history, search-affinity, competitor block). The competitor block dominates the right tail; only \texttt{has\_visitor\_history} replaces these columns with a binary indicator.}
\label{fig:missing}
\end{figure}

The \texttt{price\_usd} column contains extreme outliers, with a maximum of USD 19.7 million on the full 4.96M-row training set. A $\log(1+x)$ transform compresses the range to roughly $[0, 16.8]$, which is friendly for any later arithmetic (within-search z-scores, deltas against \texttt{prop\_log\_historical\_price}, and ratios). We do not winsorise or clip; the choice of price transformation is made on the basis of an early ablation against the within-search relativisations described in Section~\ref{subsec:feat}.

Categorical cardinality matters for encoding decisions. Small-cardinality columns (\texttt{site\_id} = 34 unique values, country IDs around 200) are passed as native LightGBM categoricals and admit straightforward one-hot encoding for XGBoost via the DMatrix categorical conversion. Large-cardinality columns (\texttt{prop\_id} with 129{,}113 unique values, \texttt{srch\_destination\_id} with 18{,}127) cannot reasonably be one-hot encoded; we pass them as LightGBM native categoricals and rely on the historical priors of Section~\ref{subsec:feat} to carry their information into XGBoost. Coverage of the test categoricals by the train slice is high (99.7\% on \texttt{prop\_id}, 97.1\% on \texttt{srch\_destination\_id}), so the global-fallback branch of the target encoder rarely fires on test.

Spearman correlations against the booking indicator (Fig.~\ref{fig:corr}) confirm two structural facts of this dataset. First, raw \texttt{price\_usd} carries essentially zero global correlation with booking ($\rho = -0.0001$): without context an absolute price tells the model nothing, because prices vary by city and currency. The same property in a different city has the same \texttt{price\_usd} but a completely different position in the user's choice set. Second, within-query relativisations of price and quality (Section~\ref{subsec:feat}) reach correlations around $0.10$ to $0.13$ with booking and $0.20$ to $0.25$ with the composite relevance grade, more than two orders of magnitude above the raw price.

\begin{figure}[t]
\centering
\includegraphics[width=0.95\linewidth]{eda_corr_booking.png}
\caption{Left, top-20 features by absolute Spearman correlation against the booking indicator on the 4.96M-row train set. Within-search relativisations and historical priors dominate the raw columns. Right, inter-correlation among the eight strongest quality features; \texttt{prop\_location\_score2} carries unique signal not redundant with price or star rating.}
\label{fig:corr}
\end{figure}

\subsection{Cleaning}

Data cleaning is deliberately conservative because LightGBM handles missing values natively and any imputation would discard the missingness signal. Five lightweight transformations are applied to both train and test. First, \texttt{date\_time} is parsed into a real timestamp so that later derivations (month, weekday, hour) can be computed without re-parsing, although they did not end up in the final feature set. Second, a binary \texttt{has\_visitor\_history} flag is added, which equals 1 if either visitor-history column is non-null. The flag has rate 5.10\% on train and 5.13\% on test, consistent. Third, \texttt{log1p\_price} is added alongside the raw \texttt{price\_usd} (Fig.~\ref{fig:price}). Fourth, three cheap competitor aggregates are computed from the row's own 24 competitor columns: \texttt{comp\_count\_quoted} (number of non-null competitor rates), \texttt{comp\_pdiff\_mean}, and \texttt{comp\_rate\_mean}. None of these aggregates cross row boundaries, so they are leakage-free by construction. Fifth, type downcasting collapses \texttt{int64} to \texttt{int32} and \texttt{float64} to \texttt{float32}, reducing the in-memory footprint from 1.72 GB to 1.16 GB on train and from 1.56 GB to 1.08 GB on test, with no measurable precision loss for tree splits.

\begin{figure}[t]
\centering
\includegraphics[width=0.85\linewidth]{cleaning_price_log_transform.png}
\caption{Price distribution before (left, clipped at USD 800 for visibility, median 122 USD) and after (right) the log1p transform (median 4.81). The original distribution has a heavy right tail up to USD 19.7 million; the log scale compresses it to a model-friendly $[0, 16.8]$ range.}
\label{fig:price}
\end{figure}

Three transformations are deliberately not applied. We do not impute competitor or visitor-history NaNs, because LightGBM splits on missingness as a separate value and that information is more useful intact. We do not winsorise \texttt{price\_usd}, because the choice of transformation depends on the downstream feature mix and the log transform is sufficient at this level. We do not encode categoricals at the cleaning stage; encoding is a feature-engineering choice that belongs in the next notebook.

\subsection{Feature engineering}\label{subsec:feat}

The locked train/val/holdout split is a group-aware 80/10/10 partition on \texttt{srch\_id} with seed 42, written to \texttt{split\_seed42.json} and reused by every later notebook. The split assigns 159{,}836 train searches (3{,}966{,}833 rows), 19{,}979 val (496{,}491 rows), and 19{,}980 holdout (495{,}023 rows). The holdout is locked: it never appears in feature engineering, training, hyperparameter selection, or fairness-rerank tuning. The val slice is used as the model-selection surface for Optuna trials, the blend weight, and the rerank lambda.

The feature set follows the published prior work on this dataset directly. The feature blocks are organised so that each is motivated by a specific reference and is added on top of the previous one only if it beats the previous on val NDCG@5 at the same anchor hyperparameters. Table~\ref{tab:fe} reports the trajectory.

\begin{table}[t]
\centering
\caption{Feature engineering ablation under a fixed LightGBM LambdaRank configuration (Bellucci~\cite{bellucci2021dmt} Table 3, 500 trees with early stopping on val NDCG@5). The hyperparameter sweep happens in Section~\ref{sec:models}.}
\label{tab:fe}
\begin{tabular}{lccc}
\toprule
Feature block & \# features & Best iter & Val NDCG@5 \\
\midrule
A: 53 raw + competitor block + categoricals & 53 & 407 & 0.39524 \\
B: + 4 within-search relativisations (Wang \& Kalousis) & 57 & 258 & 0.39884 \\
C: + 4 \texttt{prop\_id} OOF priors (Bellucci) & 61 & 263 & 0.40111 \\
D: + 4 \texttt{srch\_destination\_id} OOF priors & 65 & 378 & 0.40250 \\
E: + 4 (\texttt{prop\_id}, \texttt{srch\_destination\_id}) pair priors & 69 & 317 & 0.40258 \\
\bottomrule
\end{tabular}
\end{table}

Block B reflects the central insight of Wang and Kalousis~\cite{wang2014expedia}: ranking is an intra-query problem and absolute features matter less than their position relative to the other candidates in the same search. We add the per-search rank of \texttt{price\_usd}, the per-search z-score of \texttt{log1p\_price}, and the per-search deltas of \texttt{prop\_starrating} and \texttt{prop\_location\_score2}. These four features cost +0.00360 NDCG@5 over the raw baseline and reduce best-iteration from 407 to 258, indicating that the model converges faster when the within-query signal is exposed explicitly rather than learned implicitly.

Blocks C, D, and E implement leakage-safe historical priors following the Bellucci~\cite{bellucci2021dmt} target-encoding recipe, modified for tractability at 4.96M rows. We replace LOO encoding with 5-fold OOF encoding on \texttt{srch\_id} groups: each train row in fold $f$ receives priors aggregated over the other four folds, so its own click and book labels never enter its own prior. Validation, holdout, and test rows receive priors aggregated over the entire train slice. Smoothing follows the Laplace form $\widehat{r}_{\text{smooth}} = (\sum y + \alpha \bar y) / (n + \alpha)$ with $\alpha = 20$; unseen keys back off to the global rate. Per key we compute four statistics: log impressions, smoothed click rate, smoothed booking rate, and smoothed mean relevance. The \texttt{prop\_id} block contributes the largest single jump (+0.00227); the destination block adds a smaller +0.00139; the pair block contributes only +0.00008 in our setting, which is much less than what the autoresearch reference for this dataset would suggest. A plausible explanation is that block B already absorbs part of the within-search signal that the pair priors would otherwise carry.

Two further design choices, considered but not adopted, deserve a sentence. First, negative-row downsampling. Bellucci~\cite{bellucci2021dmt} explicitly tested downsampling negatives and reported overfitting, so they kept all rows. A small pilot at a 1:4 negatives-per-positive ratio confirms a 0.006 NDCG@5 drop on this dataset; the recipe transfers. Second, the per-property mean-position proxy from \texttt{random\_bool} = 1 rows, used by Bellucci, relies on the assumption that those impressions are position-uniform. On this dataset the position-by-\texttt{random\_bool} distribution is in fact decreasing rather than uniform, so the assumption underlying the proxy does not hold; the gain is below noise and the feature is not added.
